\documentclass[11pt]{article}
\usepackage[a4paper,margin=27mm]{geometry}
\usepackage{amsmath,amssymb,amsthm}
\newtheorem{theorem}{Theorem}
\title{An odd-prime obstruction for rational equal-weight quadrature}
\author{Research note, 30 September 2026}
\date{}
\begin{document}
\maketitle

\begin{theorem}
Let $p$ be an odd prime. Suppose rational numbers $y_1,\ldots,y_K$
form an equal-weight quadrature rule of degree at least
$\max\{p-1,3\}$ for the
uniform probability measure on $[-1,1]$. Then $K\ge 2p$.
The nodes need not belong to $[-1,1]$ and repetitions are permitted.
\end{theorem}
\begin{proof}
Suppose $K<2p$. The moment identities are
\[
 \sum_{i=1}^K y_i^s=
 \begin{cases}0,&s\text{ odd},\\K/(s+1),&s\text{ even},\end{cases}
 \qquad 1\le s\le\max\{p-1,3\}.
\]
Write $v_p$ for the $p$-adic valuation.

First suppose all $y_i$ are $p$-adically integral.
The identity at $s=p-1$ forces $p\mid K$, so $K=p$.
Reducing that identity modulo $p$, Fermat's theorem shows that the
number $q$ of unit nodes satisfies $q\equiv1\pmod p$.
Since $0\le q\le p$, exactly one node is a unit.
The first moment then cannot vanish modulo $p$, a contradiction.

Otherwise let $a=-\min_i v_p(y_i)>0$, and set $z_i=p^a y_i$.
All $z_i$ are $p$-adically integral and at least one is a unit.
For $0\le s\le p-2$ we have
\[
 \sum_{i:\,p\nmid z_i} z_i^s\equiv0\pmod p.
\]
For $s=1,\ldots,p-2$ this follows from the scaled moment identities;
their denominators $s+1$ are prime to $p$.
For $s=0$, first use $s=p-1$: its scaled right hand side
$p^{a(p-1)}K/p$ is divisible by $p$, hence the number of units is
divisible by $p$.
It lies between $1$ and $K<2p$, so it is precisely $p$.

For $b\in\mathbb F_p^*$ let $N_b$ count the unit nodes with residue
$b$. The displayed congruences give
\[
 \sum_{b\in\mathbb F_p^*}N_b b^s=0
 \quad(0\le s\le p-2).
\]
The Vandermonde matrix on the $p-1$ distinct nonzero residues is
invertible. Thus every $N_b$ is divisible by $p$.
Since their sum is $p$, all $p$ unit nodes have a common nonzero
residue $b$.

If $p=3$, for any $p$-adically integral $z\equiv b\pmod p$,
the binomial theorem gives $z^p\equiv b^p\pmod {p^2}$.
Nonunit nodes have $z_i^p\equiv0\pmod {p^2}$.
Consequently
\[
 \sum_i z_i^p\equiv p b^p\not\equiv0\pmod {p^2},
\]
contradicting the vanishing third moment.

If $p\ge5$, write the unit nodes as $z_i=b+p t_i$, and let
$T=\sum_{i:\,p\nmid z_i}t_i\pmod p$.
The scaled second and third moments are both zero modulo $p^2$:
the second is $Kp^{2a}/3$, and the third is zero.
Nonunit nodes contribute zero modulo $p^2$ to both moments.
Expanding the unit contributions gives
\[
 0\equiv p b^2+2p bT\pmod {p^2},\qquad
 0\equiv p b^3+3p b^2T\pmod {p^2}.
\]
Since $b$ is nonzero modulo $p$, these imply $b+2T=b+3T=0$
in $\mathbb F_p$, a contradiction.
\end{proof}

In particular, rational degree-three rules need at least six nodes,
and rational degree-four rules need at least ten nodes. The first bound
is sharp: the multiset $(-1,0,0,0,0,1)$ is the six-node Simpson rule.
These bounds complement the stronger asymptotic bound
$K\ge2^{\lfloor r/2\rfloor}$ proved in
\texttt{../rational\_designs/two\_adic\_lower\_bound.tex}.
Affine rescaling transfers the statement to the uniform measure on
$[0,1]$.
\end{document}
